\chapter{Mania}

\section*{Definition}
Mania is a distinct period of abnormally and persistently elevated, expansive, or irritable mood and abnormally increased goal-directed activity or energy, lasting at least 1 week (or any duration if hospitalization is required). Hypomania is a milder form lasting at least 4 consecutive days, with similar symptoms but lesser functional impairment, no psychotic features, and not requiring hospitalization. Both are core syndromes of bipolar spectrum disorders, distinguished by duration and severity.

\section*{What Happens in Your Body}
Mania is hypothesized to result from dysregulation of monoaminergic neurotransmission, particularly increased dopamine and norepinephrine activity, and altered glutamate signaling via NMDA receptors. Frontotemporal and limbic circuits (prefrontal cortex, amygdala, hippocampus, striatum) show functional and structural abnormalities. The kindling hypothesis suggests that initial mood episodes occur spontaneously after sufficient stressors; subsequent episodes become more frequent and autonomous. Circadian rhythm disruption (CLOCK gene variants, melatonin dysregulation) and mitochondrial dysfunction further contribute to mood instability.

\section*{Causes}
\begin{itemize}
  \item Primary psychiatric disorder: Bipolar I disorder (mania required), Bipolar II disorder (hypomania with major depression), Cyclothymia (chronic fluctuating hypomanic and depressive symptoms), Schizoaffective disorder (bipolar type)
  \item Secondary (medical) causes: Hyperthyroidism, Cushing disease, Cerebrovascular disease (stroke affecting right orbitofrontal or basotemporal regions), Traumatic brain injury, Multiple sclerosis (especially frontal lesions), Brain tumors (frontal, temporal, diencephalic), HIV encephalopathy, Parkinson disease medications
  \item Substance-induced: Stimulants (cocaine, amphetamines, methamphetamine, MDMA), Dopaminergic agonists (levodopa, pramipexole, ropinirole), Corticosteroids (especially at high doses), Anabolic-androgenic steroids, Antidepressant-induced (SSRIs, bupropion in bipolar people), Alcohol (intoxication or withdrawal), Hallucinogens (PCP, LSD), Cannabis (particularly high-THC strains)
  \item Sleep deprivation: Can trigger mania in susceptible individuals (bipolar vulnerability)
\end{itemize}

\section*{When to See a Doctor}
See your doctor if:
\begin{itemize}
  \item Periods of decreased need for sleep with increased energy, rapid speech, and racing thoughts lasting several days
  \item Grandiose delusions, unrealistic plans, or excessive spending sprees
  \item Reckless behavior with potential for harm (risky sexual activity, dangerous driving, poor financial decisions)
  \item Extreme irritability or agitation that is out of character and escalating
  \item Psychotic features (delusions, hallucinations) during mood elevation
  \item Any suspicion of mania in you taking antidepressants (possible unmasking of bipolar disorder)
  \item Dangerous behavior that threatens safety of self or others
\end{itemize}

\section*{Related Symptoms}
\begin{itemize}
  \item Decreased need for sleep
  \item Pressured speech
  \item Flight of ideas or racing thoughts
  \item Grandiosity
  \item Distractibility
  \item Increased goal-directed activity
  \item Psychomotor agitation
  \item Excessive involvement in pleasurable activities with high risk
  \item Poor judgment and insight
  \item Irritability or hostility when challenged
  \item Psychotic features (mood-congruent delusions, hallucinations)
  \item Catatonia
\end{itemize}
\textbf{Important:} Hypomania may be ego-syntonic and perceived positively by you, leading to treatment non-adherence; collateral history from family is essential for accurate diagnosis.

\section*{Self-Care / Home Management}
\begin{itemize}
  \item Maintain strict sleep schedule: 7--9 hours per night, as sleep deprivation is a potent trigger for mania.
  \item Avoid stimulants (caffeine, nicotine, decongestants) and all recreational drugs including alcohol.
  \item Engage a trusted family member or friend to monitor mood, spending, and behavior; grant them early warning authority.
  \item Use a mood tracking app or paper chart to log sleep, mood, energy, and medication adherence daily.
  \item Have a written relapse prevention plan with early warning signs and emergency contacts.
  \item Self-management is not appropriate during acute manic episodes; professional stabilization is required.
\end{itemize}

\section*{How Your Doctor Will Evaluate You}
\begin{itemize}
  \item History from patient and collateral sources (family is essential as insight is typically impaired) documenting mood, energy, sleep, speech, behavior, and functional impact.
  \item Screening tools: Mood Disorder Questionnaire, Young Mania Rating Scale (YMRS), Altman Self-Rating Mania Scale.
  \item Differential diagnosis: rule out medical causes (TSH, vitamin B12, toxicology screen, brain imaging if indicated) and substance-induced mood disorder.
  \item Longitudinal assessment is essential to differentiate bipolar I from bipolar II (presence/absence of mania vs. hypomania) and from borderline personality disorder (mood instability is typically hours, not days; relationship-triggered).
  \item Evaluate for current or past depressive episodes to establish bipolar spectrum diagnosis.
\end{itemize}

\section*{Treatment Options}
\begin{itemize}
  \item Acute mania: Mood stabilizers first-line (lithium --- gold standard for classic euphoric mania; divalproex or carbamazepine for mixed states or rapid cycling).
  \item Atypical antipsychotics for acute mania: olanzapine, risperidone, quetiapine, aripiprazole, ziprasidone, asenapine, or cariprazine (particularly effective for psychotic mania and severe agitation).
  \item Severe/refractory mania: combination of mood stabilizer plus atypical antipsychotic, or ECT for life-threatening or treatment-resistant cases.
  \item Avoid antidepressants as monotherapy which may worsen cycling or induce mania.
  \item Maintenance therapy: long-term lithium (reduces suicide risk, best evidence for prophylaxis), or continued mood stabilizer/atypical antipsychotic; monitor levels, renal function, and thyroid for lithium; LFTs and CBC for valproate.
  \item Psychoeducation for patient and family: early recognition of warning signs, medication adherence, sleep hygiene, and avoidance of drugs and alcohol.
  \item Cognitive-behavioral therapy and interpersonal/social rhythm therapy (IPSRT) to stabilize routines and prevent recurrence.
\end{itemize}

\clearpage
